\documentclass[a4paper,12pt]{extarticle}
\usepackage[margin=1.5cm]{geometry}
\usepackage{amsmath, amssymb}
\usepackage{array}
\usepackage{multicol}
\usepackage{helvet}
\usepackage{multirow}
\usepackage{graphicx}


\usepackage{amsmath, amsfonts, amssymb}
\usepackage{xcolor}
\usepackage{tikz}
\usepackage{enumitem}

\definecolor{numcolor}{RGB}{180, 0, 0}

\pagestyle{empty}


\renewcommand{\baselinestretch}{2}
\renewcommand{\familydefault}{\sfdefault}

\begin{document}

    \begin{tabular}{| c | c | c | c |}
        \hline
        & MATEMÁTICAS  &Fecha:$\hspace{2.5cm}$   & Calificación \\
        \hline 

        \multirow{2}{*}{\includegraphics[scale=0.15]{../../img-Logos/IesSierraCarrascoy.png}} & Nombre y apellidos \hspace{120pt} & Ex. 3º ESO   & \\ 
        &&UD2& \\ 
        \hline
    \end{tabular}
    
    \begin{tabular}{| l | }
        \hline
        NOTAS: \\ 
            1. El examen tiene que ser limpio, ordenado y sin faltas de ortografía. \\
            
            2. El examen ha de realizarse en bolígrafo negro o azul, evitando tachones en la medida \\
            de lo posible. \\ 
            
            3. Deben aparecer todas las operaciones NO vale con indicar solo el resultado.\\
            
            4. Los problemas deben contener: Datos, planteamiento y resolución, respondiendo a lo \\
            que se pregunte, NO vale con indicar un número como solución del problema.\\

            5. Persona que se vea copiando o hablando con un compañero tiene un cero en el examen\\

            6. \textbf{NO} está permitido el uso de la calculadora. \\ 
        \hline
    \end{tabular}

    \begin{enumerate}
        %potencias 
        \item \textbf{(2,5 puntos)} Simplifica aplicando las propiedades de las potencias y expresa el resultado con exponentes positivos. \textbf{Es imprescindible mostrar todo el desarrollo paso a paso; los ejercicios que solo presenten el resultado final no serán puntuados.}
            \begin{enumerate}
                \item (0,5 puntos) $(6a)^{-1 }\cdot (3a^{-2})^{-2}$
                \item (0,5 puntos) $(a^8 : a^5)^4$ 
                \item (0,5 puntos) $4^2 \cdot \left(\dfrac{3}{2}\right)^4$
                \item (1 punto)  $\dfrac{(-2)^{3} \cdot 3^{-4} \cdot 6^{2} \cdot (-3)^{5}}{(-2)^{4} \cdot 3^{3} \cdot 2^{-5} \cdot (-3)^{-2}}$
            \end{enumerate}

        %raices
        \item \textbf{(3 puntos)} Realiza las siguientes operaciones con radicales \textbf{sin calculadora} y simplifica al máximo el resultado \textbf{(es imprescindible mostrar todos los pasos para puntuar)}:
        \begin{enumerate}[label=\alph*)]
            \item \textbf{Extrae todos los factores posibles:} $\sqrt[3]{250 \cdot x^{7} \cdot y^{11} \cdot z^{3}}$
            \item \textbf{Efectúa y simplifica:} $\sqrt[3]{3} \cdot \sqrt[3]{1125}$
            \item \textbf{Efectúa la suma y resta de radicales:} $3\sqrt{48} - 2\sqrt{75} + \dfrac{1}{2}\sqrt{12} - \sqrt{27}$
        \end{enumerate}
        

        %notación cientifica
        \item \textbf{(1,5 puntos)} Realiza la siguiente operación combinada \textbf{sin calculadora} y expresa el resultado final en notación científica (es imprescindible mostrar todos los pasos).
        \begin{enumerate}
            \item $(1,5 \cdot 10^{-4}) \cdot (4 \cdot 10^{9}) - 2,5 \cdot 10^{5}$
            \item $\dfrac{(3,6 \cdot 10^{8} - 1,1 \cdot 10^{7}) \cdot 2 \cdot 10^{-3}}{5 \cdot 10^{-4}}$
        \end{enumerate}
        
        %problema de porcentajes
         \item \textbf{(1,5 punto)} Para desbloquear un aspecto legendario en su videojuego favorito, Alex ya ha reunido $480$ monedas, que representan el $80\%$ de su coste total. ¿Cuántas monedas cuesta el aspecto legendario? \textbf{(Es imprescindible mostrar el planteamiento y las operaciones paso a paso; sin desarrollo no se puntuará).}
        \item \textbf{(1,5 punto)} La entrada a un parque de atracciones costaba inicialmente $40\text{ €}$. En temporada alta su precio subió un $25\%$, pero después lanzaron una oferta con un descuento del $20\%$ sobre el nuevo precio. ¿Cuánto cuesta la entrada finalmente? \textbf{(Es imprescindible mostrar todo el desarrollo paso a paso; sin operaciones no se puntuará).}
    \end{enumerate}


\end{document}